\section{Introduction}\label{sec:introduction}

Digital parties often claim to make politics more responsive by lowering the cost of participation. The central question is whether this voice changes what parties do. This paper studies that question in the case of Beppe Grillo's blog and the Five Star Movement (M5S), where a party-owned digital platform was central to movement identity, mobilization, agenda formation, and the relationship between leaders and supporters.

The project asks whether grievances, priorities, local information, and policy proposals expressed in blog comments later shaped Grillo/M5S posts, campaign materials, programs, or parliamentary outputs. The analysis is planned as a longitudinal text design linking supporter input, party communication, formal political outputs, and external agenda controls.

The main findings are TBD. Until data recovery and validation are complete, the paper should be framed as a research design and measurement project rather than as evidence of responsiveness.
